\documentclass[11pt,a4paper]{article}
\usepackage{templates/modeling-notes}
\usepackage{needspace}
\hypersetup{pdftitle={Chapter 4 -- Degeneracy, duality, and dual simplex}}
\title{\textbf{Mathematical Modeling}\\[2mm]
\large Chapter 4: Degeneracy, Duality, and Dual Simplex}
\author{Personal study notes}
\date{}
\begin{document}
\maketitle
{\small\tableofcontents}
\clearpage

\section{A common notation for both simplex methods}\label{sec:notation}
Consider the standard maximization problem
\[
  \max c^Tx,\qquad Ax=b,\qquad x\geq0,
  \qquad A\in\R^{m\times n},\quad \operatorname{rank}A=m.
\]
An ordered basis $B=(j_1,\ldots,j_m)$ selects a nonsingular $A_B$; $N$
contains the nonbasic indices. Define
\begin{equation}\label{eq:canonical}
  \bar A=A_B^{-1}A,\qquad \bar b=A_B^{-1}b,\qquad
  q=\bar A^Tc_B-c,\qquad z_0=c_B^T\bar b.
\end{equation}
The basic solution has $x_N=0$ and $x_B=\bar b$. Its equations and objective are
\begin{equation}\label{eq:dictionary}
  x_B=\bar b-\bar A_Nx_N,\qquad
  z(x_N)=z_0-\sum_{j\in N}q_jx_j.
\end{equation}
Throughout these notes, $q_j=z_j-c_j$ is an \emph{optimality indicator}.
The more usual reduced cost is $\rho_j=-q_j$; never mix their signs.
Row $r$ corresponds to variable $x_{j_r}$, not necessarily $x_r$.

\begin{center}
  \begin{tabular}{lll}\toprule
    Condition & Meaning & Simplex role\\\midrule
    $\bar b\geq0$ & Primal feasible basic solution & Primal method preserves it\\
    $q\geq0$ & Dual feasible basis & Dual method preserves it\\
    Both conditions & Optimal primal and dual pair & Stop with a certificate\\
    $\bar b_i=0$ for some $i$ & Degenerate feasible basis & Zero steps may occur\\\bottomrule
  \end{tabular}
\end{center}
A primal pivot chooses $k\in N$ with $q_k<0$. For $u=\bar A_{:k}$,
$x_k=\theta$ gives $x_B=\bar b-\theta u$ and objective change $-q_k\theta$.
If some $u_i>0$, the largest feasible step is
\[
  \theta^*=\min_{i:u_i>0}\frac{\bar b_i}{u_i}.
\]
If $u\leq0$, the objective is unbounded above along this feasible ray.
Both methods use the same Gauss--Jordan pivot operation. If $p=\bar a_{rk}\ne0$,
\begin{align*}
  \bar a'_{rj}&=\bar a_{rj}/p,& \bar b'_r&=\bar b_r/p,\\
  \bar a'_{ij}&=\bar a_{ij}-\bar a_{ik}\bar a_{rj}/p,
  &\bar b'_i&=\bar b_i-\bar a_{ik}\bar b_r/p\quad(i\ne r),\\
  q'_j&=q_j-q_k\bar a_{rj}/p,& z'_0&=z_0-q_k\bar b_r/p.
\end{align*}

\section{Degeneracy: changing a basis without changing the point}\label{sec:degeneracy}
A basic feasible solution is degenerate when at least one \emph{basic}
coordinate is zero. Having nonbasic zeros is normal and says nothing about
degeneracy. Two mechanisms deserve separate attention.
\begin{center}
  \begin{tabular}{>{\raggedright\arraybackslash}p{3.2cm}>{\raggedright\arraybackslash}p{5.3cm}>{\raggedright\arraybackslash}p{6cm}}\toprule
    Ratio test & Next point & Objective for $q_k<0$\\\midrule
    $\theta^*=0$ & Same full vector $x$; a different basis & No change\\
    $\theta^*>0$, unique minimum & A different point & Strict improvement\\
    $\theta^*>0$, tied minimum & A different, degenerate point: another basic variable also reaches zero & Strict improvement\\\bottomrule
  \end{tabular}
\end{center}
For a tied positive ratio, $\bar b_i-\theta^*u_i=0$ for every tied row.
Only one variable leaves the basis; the others become zero basic variables.
Thus degeneracy may be reached by a strictly improving step.
A zero step need not prove that the current point is suboptimal: a degenerate
optimal point can have a basis with some negative indicators.

\begin{figure}[H]
  \centering
  \begin{tikzpicture}[every node/.style={font=\small},
      box/.style={draw=modelblue,rounded corners,align=center,minimum width=3.5cm,minimum height=1cm}]
    \node[box] (b1) at (0,0) {Basis $B$\\one basic zero};
    \node[box] (b2) at (5.8,0) {Basis $B'$\\one basic zero};
    \node[model point,label=below:{same feasible point $x$}] (x) at (2.9,-1.5) {};
    \draw[model flow] (b1)--node[above] {$\theta^*=0$} (b2);
    \draw[dashed,modelblue] (b1)--(x);
    \draw[dashed,modelblue] (b2)--(x);
    \node at (2.9,1) {Basis exchange $\ne$ geometric movement};
  \end{tikzpicture}
  \caption{A degenerate pivot can change the equations used to describe a vertex
    without moving away from that vertex.}
\end{figure}
There are finitely many bases. If every pivot strictly improves the objective,
no basis can recur. Zero steps remove this argument. With deterministic choices,
a return to the same ordered basis repeats the same tableau and can initiate
an endless cycle. Repeated points alone do not establish a cycle: several
successive bases can describe the same point before the algorithm escapes.

\subsection{A six-pivot cycling example}\label{sec:cycling}
Consider seven nonnegative variables and the problem
\begin{equation}\label{eq:beale}
  \begin{aligned}
    \max\quad &\tfrac34x_4-20x_5+\tfrac12x_6-6x_7,\\
    &x_1+\tfrac14x_4-8x_5-x_6+9x_7=0,\\
    &x_2+\tfrac12x_4-12x_5-\tfrac12x_6+3x_7=0,\\
    &x_3+x_6=1,\qquad x_1,\ldots,x_7\geq0.
  \end{aligned}
\end{equation}
Start at $B=(1,2,3)$, so $x=(0,0,1,0,0,0,0)^T$ and $z_0=0$.
Use Dantzig's entering rule (most negative $q_j$) and resolve a tied ratio
by taking the first row in the current tableau. The successive tableaus below
show the complete cycle. The bottom row is $q^T$ with $z_0$ on the right;
basic columns are unit columns.
\Needspace{8\baselineskip}
\noindent\textbf{Basis $B=(1,2,3)$.}
\begin{center}
  \begin{tabular}{c|rrrrrrr|r}\toprule
    Basic & $x_1$ & $x_2$ & $x_3$ & $x_4$ & $x_5$ & $x_6$ & $x_7$ & RHS\\\midrule
    $x_1$ & 1 & 0 & 0 & $1/4$ & -8 & -1 & 9 & 0\\
    $x_2$ & 0 & 1 & 0 & $1/2$ & -12 & $-1/2$ & 3 & 0\\
    $x_3$ & 0 & 0 & 1 & 0 & 0 & 1 & 0 & 1\\
    \midrule
    $q^T$ & 0 & 0 & 0 & $-3/4$ & 20 & $-1/2$ & 6 & 0\\\bottomrule
  \end{tabular}
\end{center}

\Needspace{8\baselineskip}
\noindent\textbf{Basis $B=(4,2,3)$.}
\begin{center}
  \begin{tabular}{c|rrrrrrr|r}\toprule
    Basic & $x_1$ & $x_2$ & $x_3$ & $x_4$ & $x_5$ & $x_6$ & $x_7$ & RHS\\\midrule
    $x_4$ & 4 & 0 & 0 & 1 & -32 & -4 & 36 & 0\\
    $x_2$ & -2 & 1 & 0 & 0 & 4 & $3/2$ & -15 & 0\\
    $x_3$ & 0 & 0 & 1 & 0 & 0 & 1 & 0 & 1\\
    \midrule
    $q^T$ & 3 & 0 & 0 & 0 & -4 & $-7/2$ & 33 & 0\\\bottomrule
  \end{tabular}
\end{center}

\Needspace{8\baselineskip}
\noindent\textbf{Basis $B=(4,5,3)$.}
\begin{center}
  \begin{tabular}{c|rrrrrrr|r}\toprule
    Basic & $x_1$ & $x_2$ & $x_3$ & $x_4$ & $x_5$ & $x_6$ & $x_7$ & RHS\\\midrule
    $x_4$ & -12 & 8 & 0 & 1 & 0 & 8 & -84 & 0\\
    $x_5$ & $-1/2$ & $1/4$ & 0 & 0 & 1 & $3/8$ & $-15/4$ & 0\\
    $x_3$ & 0 & 0 & 1 & 0 & 0 & 1 & 0 & 1\\
    \midrule
    $q^T$ & 1 & 1 & 0 & 0 & 0 & -2 & 18 & 0\\\bottomrule
  \end{tabular}
\end{center}

\Needspace{8\baselineskip}
\noindent\textbf{Basis $B=(6,5,3)$.}
\begin{center}
  \begin{tabular}{c|rrrrrrr|r}\toprule
    Basic & $x_1$ & $x_2$ & $x_3$ & $x_4$ & $x_5$ & $x_6$ & $x_7$ & RHS\\\midrule
    $x_6$ & $-3/2$ & 1 & 0 & $1/8$ & 0 & 1 & $-21/2$ & 0\\
    $x_5$ & $1/16$ & $-1/8$ & 0 & $-3/64$ & 1 & 0 & $3/16$ & 0\\
    $x_3$ & $3/2$ & -1 & 1 & $-1/8$ & 0 & 0 & $21/2$ & 1\\
    \midrule
    $q^T$ & -2 & 3 & 0 & $1/4$ & 0 & 0 & -3 & 0\\\bottomrule
  \end{tabular}
\end{center}

\Needspace{8\baselineskip}
\noindent\textbf{Basis $B=(6,7,3)$.}
\begin{center}
  \begin{tabular}{c|rrrrrrr|r}\toprule
    Basic & $x_1$ & $x_2$ & $x_3$ & $x_4$ & $x_5$ & $x_6$ & $x_7$ & RHS\\\midrule
    $x_6$ & 2 & -6 & 0 & $-5/2$ & 56 & 1 & 0 & 0\\
    $x_7$ & $1/3$ & $-2/3$ & 0 & $-1/4$ & $16/3$ & 0 & 1 & 0\\
    $x_3$ & -2 & 6 & 1 & $5/2$ & -56 & 0 & 0 & 1\\
    \midrule
    $q^T$ & -1 & 1 & 0 & $-1/2$ & 16 & 0 & 0 & 0\\\bottomrule
  \end{tabular}
\end{center}

\Needspace{8\baselineskip}
\noindent\textbf{Basis $B=(1,7,3)$.}
\begin{center}
  \begin{tabular}{c|rrrrrrr|r}\toprule
    Basic & $x_1$ & $x_2$ & $x_3$ & $x_4$ & $x_5$ & $x_6$ & $x_7$ & RHS\\\midrule
    $x_1$ & 1 & -3 & 0 & $-5/4$ & 28 & $1/2$ & 0 & 0\\
    $x_7$ & 0 & $1/3$ & 0 & $1/6$ & -4 & $-1/6$ & 1 & 0\\
    $x_3$ & 0 & 0 & 1 & 0 & 0 & 1 & 0 & 1\\
    \midrule
    $q^T$ & 0 & -2 & 0 & $-7/4$ & 44 & $1/2$ & 0 & 0\\\bottomrule
  \end{tabular}
\end{center}

\Needspace{8\baselineskip}
\noindent\textbf{Basis $B=(1,2,3)$.}
\begin{center}
  \begin{tabular}{c|rrrrrrr|r}\toprule
    Basic & $x_1$ & $x_2$ & $x_3$ & $x_4$ & $x_5$ & $x_6$ & $x_7$ & RHS\\\midrule
    $x_1$ & 1 & 0 & 0 & $1/4$ & -8 & -1 & 9 & 0\\
    $x_2$ & 0 & 1 & 0 & $1/2$ & -12 & $-1/2$ & 3 & 0\\
    $x_3$ & 0 & 0 & 1 & 0 & 0 & 1 & 0 & 1\\
    \midrule
    $q^T$ & 0 & 0 & 0 & $-3/4$ & 20 & $-1/2$ & 6 & 0\\\bottomrule
  \end{tabular}
\end{center}

Every step has length zero. The first entry is $x_4$, with ratios $0/(1/4)$
and $0/(1/2)$; choosing row 1 makes $x_1$ leave. At the fourth tableau,
$q_1=-2$ and $q_7=-3$, so Dantzig's rule chooses $x_7$.
After six pivots, both the basis and the full tableau are back at their start.

\begin{figure}[H]
  \centering
  \begin{tikzpicture}[every node/.style={font=\small},
      box/.style={draw=modelblue,rounded corners,minimum width=2.3cm,minimum height=.7cm}]
    \node[box] (a) at (0,0) {$(1,2,3)$};
    \node[box] (b) at (4,0) {$(4,2,3)$};
    \node[box] (c) at (8,0) {$(4,5,3)$};
    \node[box] (d) at (8,-1.6) {$(6,5,3)$};
    \node[box] (e) at (4,-1.6) {$(6,7,3)$};
    \node[box] (f) at (0,-1.6) {$(1,7,3)$};
    \draw[model flow] (a)--node[above] {$4\gets1$} (b);
    \draw[model flow] (b)--node[above] {$5\gets2$} (c);
    \draw[model flow] (c)--node[right] {$6\gets4$} (d);
    \draw[model flow] (d)--node[below] {$7\gets5$} (e);
    \draw[model flow] (e)--node[below] {$1\gets6$} (f);
    \draw[model flow] (f)--node[left] {$2\gets7$} (a);
    \node at (4,-2.8) {Every pivot: $\theta^*=0$, $x=(0,0,1,0,0,0,0)^T$, $z_0=0$};
  \end{tikzpicture}
  \caption{A cycle in the space of bases. Labels mean ``entering index $\gets$
    leaving index''; there is no corresponding tour of distinct vertices.}
\end{figure}

\subsection{Bland's rule and its complete path}\label{sec:bland}
Bland's anti-cycling rule changes \emph{both} choices:
\begin{enumerate}[leftmargin=*,itemsep=2pt]
  \item Choose the smallest original nonbasic index $k$ with $q_k<0$.
  \item Apply the ordinary minimum-ratio test; among ties, choose the smallest
        original index $j_r$ of the basic variable that would leave.
\end{enumerate}
In exact arithmetic this rule prevents cycling and yields finite termination
of primal simplex. The claim is not a promise about speed or floating-point
roundoff. A rule changing only the entering choice is not the full Bland rule.

\begin{center}
  \begin{tabular}{clcccl}\toprule
    Step & Ordered basis & Enter & Leave & $\theta^*$ & $z_0$\\\midrule
    0 & $(1,2,3)$ & 4 & 1 & 0 & 0\\
    1 & $(4,2,3)$ & 5 & 2 & 0 & 0\\
    2 & $(4,5,3)$ & 6 & 4 & 0 & 0\\
    3 & $(6,5,3)$ & 1 & 5 & 0 & 0\\
    4 & $(6,1,3)$ & 2 & 3 & $1/2$ & 0\\
    5 & $(6,1,2)$ & 4 & 2 & 1 & $1/2$\\
    6 & $(6,1,4)$ & --- & --- & --- & $5/4$\\\bottomrule
  \end{tabular}
\end{center}
At step 3, Bland chooses $x_1$ rather than $x_7$, breaking the previous cycle.
For that column the only positive transformed coefficient is $1/16$ in the
$x_5$ row, so $x_5$ leaves at zero step. The remaining pivots move to
\[
  x^*=(3/4,0,0,1,0,1,0)^T,\qquad z^*=5/4.
\]
The final indicators are $q=(0,3/2,5/4,0,2,0,21/2)^T\geq0$.
Equation~\eqref{eq:dictionary} therefore gives the exact certificate
\begin{equation}\label{eq:beale-certificate}
  z=\tfrac54-\tfrac32x_2-\tfrac54x_3-2x_5-\tfrac{21}{2}x_7\leq\tfrac54.
\end{equation}
All four nonbasic indicators are strictly positive. Equality forces their
variables to zero, after which the basis equations fix the remaining variables;
the optimum is unique.
\paragraph{Practice prompt.}
\foreignlanguage{polish}{Rozwiązać przykład Beale, stosując regułę Bland'a.}
The path above supplies a hand solution; the MATLAB companion checks its choices,
steps, and certificate.

\Needspace{14\baselineskip}
\section{Perturbation and lexicographic tie-breaking}\label{sec:perturbation}
Degeneracy also has a column-space interpretation. For a fixed basis,
\[
  b=\sum_{i=1}^m a_{j_i}\bar b_i.
\]
A feasible basis writes $b$ as a nonnegative combination of its columns.
A degenerate basis places $b$ on the boundary of that simplicial cone: at least
one coefficient is zero. A small symbolic perturbation can place it in the
interior of a nearby cone.
\begin{figure}[H]
  \centering
  \begin{tikzpicture}[every node/.style={font=\small},scale=1.05]
    \fill[modelblue!10] (0,0)--(4.2,0)--(3,3)--cycle;
    \draw[model axis] (0,0)--(4.8,0) node[right] {first coordinate};
    \draw[model axis] (0,0)--(0,3.5) node[above] {second coordinate};
    \draw[model flow] (0,0)--(4.2,0) node[below] {$a_{j_1}$};
    \draw[model flow] (0,0)--(3,3) node[above] {$a_{j_2}$};
    \draw[modelorange,thick,-{Stealth}] (0,0)--(2.5,0);
    \node[model point,label=below:{$b$: boundary}] at (2.5,0) {};
    \draw[modelorange,thick,-{Stealth}] (2.5,.05)--(2.7,.8);
    \node[model point,label=right:{$b(\varepsilon)$: positive coefficients}] at (2.7,.8) {};
  \end{tikzpicture}
  \caption{A two-dimensional illustration in the space of right-hand sides.
    The decision vector and the right-hand-side vector belong to different spaces.}
\end{figure}

\subsection{Symbolic infinitesimals}\label{sec:lex}
Relabel columns so the initial feasible basis is $B_0=(1,\ldots,m)$.
Consider the formal perturbation
\begin{equation}\label{eq:perturb}
  b(\varepsilon)=b+\sum_{j=1}^n\varepsilon^j a_j,
  \qquad 0<\varepsilon\ll1.
\end{equation}
Here $\varepsilon^2$ is much smaller than $\varepsilon$, and so on.
At the initial basis, the first $m$ transformed columns are the identity.
Consequently each zero basic coordinate gets a positive leading perturbation,
so this basis stays feasible for sufficiently small positive $\varepsilon$.
Without an initial ordering or a perturbation adapted to $B_0$, feasibility
of the starting basis must be checked; arbitrary RHS noise is not this rule.

For any basis,
\[
  \bar b_i(\varepsilon)=\bar b_i+\sum_{j=1}^n\varepsilon^j\bar a_{ij}.
\]
Every such polynomial is nonzero as a polynomial: a basic unit column supplies
a nonzero coefficient. There are finitely many bases and rows, hence only
finitely many positive roots. Choose $\varepsilon$ below their smallest positive
root, if there is one. Then no basic coordinate equals zero in the perturbed
problem. Feasible pivots with $q_k<0$ strictly improve its objective and cannot
cycle. This is an exact-arithmetic argument, not a recommended numerical choice
of a tiny floating-point constant.

When $\bar a_{ik}>0$, compare the coefficient tuples
\begin{equation}\label{eq:lex-tuple}
  L_i=\left(\frac{\bar b_i}{\bar a_{ik}},
  \frac{\bar a_{i1}}{\bar a_{ik}},\ldots,
  \frac{\bar a_{in}}{\bar a_{ik}}\right)
\end{equation}
\emph{lexicographically}: compare the first coordinate, then the first coordinate
where the tuples differ. Choose the smallest tuple. This implements the
perturbed ratio test without assigning a numerical value to $\varepsilon$.
Distinct normalized rows cannot coincide, since $\bar A$ has full row rank.
An equivalent shorter perturbation uses
$b+A_{B_0}(\varepsilon,\ldots,\varepsilon^m)^T$ and keeps the transformed
$B_0$ columns as its tie-break coefficients.

\Needspace{6\baselineskip}
\paragraph{First tie in the cycling example.}
For $k=4$, the ordinary ratio is zero in both first rows. Their tuples begin
\[
  L_1=(0,4,0,0,\ldots),\qquad L_2=(0,0,2,0,\ldots).
\]
Since $0<4$ in the first differing position, row 2 wins: $x_2$ leaves.
This gives the following two noninitial tableaus.
\Needspace{8\baselineskip}
\noindent\textbf{Basis $B=(1,4,3)$.}
\begin{center}
  \begin{tabular}{c|rrrrrrr|r}\toprule
    Basic & $x_1$ & $x_2$ & $x_3$ & $x_4$ & $x_5$ & $x_6$ & $x_7$ & RHS\\\midrule
    $x_1$ & 1 & $-1/2$ & 0 & 0 & -2 & $-3/4$ & $15/2$ & 0\\
    $x_4$ & 0 & 2 & 0 & 1 & -24 & -1 & 6 & 0\\
    $x_3$ & 0 & 0 & 1 & 0 & 0 & 1 & 0 & 1\\
    \midrule
    $q^T$ & 0 & $3/2$ & 0 & 0 & 2 & $-5/4$ & $21/2$ & 0\\\bottomrule
  \end{tabular}
\end{center}

\Needspace{8\baselineskip}
\noindent\textbf{Basis $B=(1,4,6)$.}
\begin{center}
  \begin{tabular}{c|rrrrrrr|r}\toprule
    Basic & $x_1$ & $x_2$ & $x_3$ & $x_4$ & $x_5$ & $x_6$ & $x_7$ & RHS\\\midrule
    $x_1$ & 1 & $-1/2$ & $3/4$ & 0 & -2 & 0 & $15/2$ & $3/4$\\
    $x_4$ & 0 & 2 & 1 & 1 & -24 & 0 & 6 & 1\\
    $x_6$ & 0 & 0 & 1 & 0 & 0 & 1 & 0 & 1\\
    \midrule
    $q^T$ & 0 & $3/2$ & $5/4$ & 0 & 2 & 0 & $21/2$ & $5/4$\\\bottomrule
  \end{tabular}
\end{center}

After this zero step, $x_6$ enters. The only positive coefficient in that
column is 1 in the $x_3$ row, so $x_3$ leaves with step 1. The optimizer is
again $(3/4,0,0,1,0,1,0)^T$ and $z=5/4$.
The perturbed objectives strictly improve even when the unperturbed step is
zero. As $\varepsilon\to0$, an optimal certifying basis of the perturbed
problem gives a feasible limit with unchanged nonnegative indicators, hence an
optimum of the original problem.

\begin{center}
  \begin{tabular}{llll}\toprule
    Rule on this example & Pivots & Behavior & Optimum\\\midrule
    Dantzig + first-row ties & 6, then repeat & Cycles at the start & Not reached\\
    Full Bland & 6 & Four zero steps, two moves & $5/4$\\
    Dantzig + lexicographic ties & 2 & One zero step, one move & $5/4$\\\bottomrule
  \end{tabular}
\end{center}
The counts compare this particular example, not universal relative performance.

\Needspace{23\baselineskip}
\section{Complexity and the route to an optimum}\label{sec:complexity}
Avoiding cycles establishes termination, but does not establish a polynomial
number of pivots. A family of distorted cubes exposes the difference.
For $p=5$, the three-dimensional member is
\[
  \begin{aligned}
    \max\quad &4x_1+2x_2+x_3,\\
    &x_1\leq5,\\
    &4x_1+x_2\leq25,\\
    &8x_1+4x_2+x_3\leq125,\qquad x\geq0.
  \end{aligned}
\]
Starting at the origin, Dantzig's rule visits all eight vertices:
\begin{center}
  \begin{tabular}{clrr}\toprule
    Step & $(x_1,x_2,x_3)$ & $4x_1+2x_2+x_3$ & Entering index next\\\midrule
    0 & $(0,0,0)$ & 0 & 1\\
    1 & $(5,0,0)$ & 20 & 2\\
    2 & $(5,5,0)$ & 30 & 4\\
    3 & $(0,25,0)$ & 50 & 3\\
    4 & $(0,25,25)$ & 75 & 1\\
    5 & $(5,5,65)$ & 95 & 5\\
    6 & $(5,0,85)$ & 105 & 4\\
    7 & $(0,0,125)$ & 125 & ---\\\bottomrule
  \end{tabular}
\end{center}
Indices 4--6 are the slack variables. The last constraint yields
\[
  4x_1+2x_2+x_3\leq125-4x_1-2x_2\leq125,
\]
with equality only at $(0,0,125)$.
The analogous $n$-variable family maximizes
$\sum_{j=1}^n2^{n-j}x_j$, with rows
\[
  x_i+\sum_{j<i}2^{i-j+1}x_j\leq p^i,\quad i=1,\ldots,n,
  \qquad x\geq0.
\]
For the usual choice $p=5$, this rule can require $2^n-1$ pivots. This is a
worst-case statement about a specified pivot rule and family. It does not say
that every simplex rule follows this path, or that simplex is slow on every LP.

\begin{figure}[H]
  \centering
  \begin{tikzpicture}[x=.8cm,y=.13cm,every node/.style={font=\small}]
    \filldraw[model region] (0,0)--(5,0)--(5,5)--(0,25)--cycle;
    \draw[model axis] (0,0)--(6,0) node[right] {$x_1$};
    \draw[model axis] (0,0)--(0,29) node[above] {$x_2$};
    \draw[modelorange,very thick,-{Stealth}] (0,0)--(5,0);
    \draw[modelorange,very thick,-{Stealth}] (5,0)--(5,5);
    \draw[modelorange,very thick,-{Stealth}] (5,5)--(0,25);
    \foreach \x/\y in {0/0,5/0,5/5,0/25} \node[model point] at (\x,\y) {};
    \node[below left] at (0,0) {$(0,0)$};
    \node[below] at (5,0) {$(5,0)$};
    \node[right] at (5,5) {$(5,5)$};
    \node[left] at (0,25) {$(0,25)$};
    \node[align=left] at (10,14) {Two-dimensional member\\$\max(2x_1+x_2)$\\$x_1\leq5$, $4x_1+x_2\leq25$};
  \end{tikzpicture}
  \caption{The distorted square already makes Dantzig's rule visit every vertex.
    The higher-dimensional family extends this long edge-following route.}
\end{figure}
For rational input, polynomial-time complexity is measured against the
\emph{bit length} of all input coefficients, not just $m+n$.
Ellipsoid and suitable interior-point algorithms have polynomial-time guarantees.
Simplex is widely effective in practice, but no universal ``average linear in
$m+n$'' bound follows without a stated probability model and analysis.
Floating-point conditioning and sparse linear algebra also affect real runtime.

\begin{figure}[H]
  \centering
  \begin{tikzpicture}[every node/.style={font=\small}]
    \foreach \shift in {0,6.4} {
        \begin{scope}[xshift=\shift cm]
          \filldraw[model region] (0,0)--(3,0)--(3,1)--(1.2,3)--(0,2.5)--cycle;
          \node[model point,label=above:{$x^*$}] at (1.2,3) {};
        \end{scope}
      }
    \draw[modelorange,very thick,-{Stealth}] (0,0)--(3,0)--(3,1)--(1.2,3);
    \node at (1.5,-.5) {Simplex: basis exchanges};
    \draw[modelorange,very thick,-{Stealth}] (7.1,.6)--(7.7,1.25)--(7.9,2.05)--(7.65,2.65);
    \foreach \x/\y in {7.1/.6,7.7/1.25,7.9/2.05,7.65/2.65}
    \fill[modelorange] (\x,\y) circle (1.6pt);
    \node at (7.9,-.5) {Interior-point: interior iterates};
  \end{tikzpicture}
  \caption{Schematic algorithmic routes, not computed iterates. An interior-point
    method approaches an optimal face; that face need not be a single vertex.
    Crossover can subsequently select a basic solution.}
\end{figure}

\section{Constructing the dual problem}\label{sec:dual-rules}
\subsection{Bounds from weighted constraints}
For the inequality-form primal
\begin{equation}\label{eq:primal-dual}
  \text{(P)}\quad\max c^Tx:\ Ax\leq b,\ x\geq0,
  \qquad
  \text{(D)}\quad\min b^Ty:\ A^Ty\geq c,\ y\geq0,
\end{equation}
think of $y_i$ as the weight or price of resource constraint $i$.
Nonnegative weights preserve its direction. If $A^Ty\geq c$, each product's
weighted resource use covers its objective coefficient, so $b^Ty$ bounds the
value achievable with available resources. Minimizing chooses the tightest bound.

The standard equality-form primal has dual
\begin{equation}\label{eq:standard-dual}
  \max c^Tx:\ Ax=b,\ x\geq0
  \quad\longleftrightarrow\quad
  \min b^Ty:\ A^Ty\geq c,\ y\in\R^m.
\end{equation}
Equality constraints allow weights of either sign. Requiring $y\geq0$ here
would generally change the dual problem.

\begin{figure}[H]
  \centering
  \begin{tikzpicture}[every node/.style={font=\small},
      box/.style={draw=modelblue,align=center,minimum width=3.4cm,minimum height=.8cm}]
    \node[box] (pc) at (0,0) {$m$ primal constraints};
    \node[box] (dv) at (7,0) {$m$ dual variables};
    \node[box] (pv) at (0,-1.4) {$n$ primal variables};
    \node[box] (dc) at (7,-1.4) {$n$ dual constraints};
    \draw[model flow] (pc)--node[above] {one weight per row} (dv);
    \draw[model flow] (pv)--node[above] {one test per column} (dc);
    \node at (3.5,-2.4) {$A\ (m\times n)\ \longmapsto\ A^T\ (n\times m)$};
  \end{tikzpicture}
  \caption{Dualization interchanges rows and columns, $b$ and $c$, and max and min.
    The number of constraints changes places with the number of variables.}
\end{figure}

\subsection{A sign table for mixed forms}
Write the primal as a \emph{maximization}. Then the dual is a minimization,
with objective $b^Ty$, and use the following two independent tables.
\begin{center}
  \begin{tabular}{ll@{\qquad}ll}\toprule
    Primal row & Dual variable & Primal variable & Dual column condition\\\midrule
    $a_i^Tx\leq b_i$ & $y_i\geq0$ & $x_j\geq0$ & $a_j^Ty\geq c_j$\\
    $a_i^Tx\geq b_i$ & $y_i\leq0$ & $x_j\leq0$ & $a_j^Ty\leq c_j$\\
    $a_i^Tx=b_i$ & $y_i$ free & $x_j$ free & $a_j^Ty=c_j$\\\bottomrule
  \end{tabular}
\end{center}
Here $a_i^T$ denotes row $i$, while $a_j$ denotes column $j$.
For a primal \emph{minimization}, reverse the inequality directions in the
column conditions and the sign directions for inequality-row weights:
\begin{center}
  \begin{tabular}{ll@{\qquad}ll}\toprule
    Primal row & Dual variable & Primal variable & Dual column condition\\\midrule
    $a_i^Tx\geq b_i$ & $y_i\geq0$ & $x_j\geq0$ & $a_j^Ty\leq c_j$\\
    $a_i^Tx\leq b_i$ & $y_i\leq0$ & $x_j\leq0$ & $a_j^Ty\geq c_j$\\
    $a_i^Tx=b_i$ & $y_i$ free & $x_j$ free & $a_j^Ty=c_j$\\\bottomrule
  \end{tabular}
\end{center}
A quick derivation for max is to form
\[
  L(x,y)=c^Tx+y^T(b-Ax)=b^Ty+(c-A^Ty)^Tx.
\]
The row sign rules ensure $L(x,y)\geq c^Tx$ on primal-feasible $x$.
The column rules ensure its supremum over the variable domains is $b^Ty$:
a positive coefficient on a nonnegative variable would otherwise give $+\infty$;
a free variable requires a zero coefficient. Thus the signs follow from an
upper-bound construction, rather than memorization alone.
Dualizing again restores the original problem, up to equivalent sign conventions.

\subsection{Worked mixed-sign example}\label{sec:mixed}
\[
  \begin{array}{ll}
    \text{Primal:}\ \max &8x_1+3x_2-2x_3\\
    \text{subject to}&x_1-6x_2+x_3\geq2,\\
    &5x_1+7x_2-2x_3=-4,\\
    &x_1\leq0,\quad x_2\geq0,\quad x_3\in\R.
  \end{array}
\]
The first row gives $y_1\leq0$; the second makes $y_2$ free. Read one column
at a time to obtain
\[
  \begin{array}{ll}
    \text{Dual:}\ \min &2y_1-4y_2\\
    \text{subject to}&y_1+5y_2\leq8,\\
    &-6y_1+7y_2\geq3,\\
    &y_1-2y_2=-2,\\
    &y_1\leq0,\quad y_2\in\R.
  \end{array}
\]
As a check, $y=(0,1)$ is dual feasible. The primal point $x=(0,0,2)$
is feasible and gives $c^Tx=-4=b^Ty$. Equal feasible values certify that
both points are optimal, as the next section proves.

\section{Weak duality, strong duality, and possible outcomes}\label{sec:duality}
\begin{theorem}[Weak duality]
  For the pair in \eqref{eq:primal-dual}, every primal-feasible $x$ and
  dual-feasible $y$ satisfy $c^Tx\leq b^Ty$.
\end{theorem}
\begin{proof}
  Using nonnegativity in the two separate comparisons,
  \[
    c^Tx\leq x^TA^Ty=y^TAx\leq y^Tb.
  \]
\end{proof}
\begin{corollary}[Equal feasible values certify optimality]
  If feasible $x,y$ satisfy $c^Tx=b^Ty$, both are optimal.
\end{corollary}
\begin{proof}
  Every feasible primal value is at most $b^Ty=c^Tx$, and every feasible dual
  value is at least $c^Tx=b^Ty$.
\end{proof}
This certificate uses only weak duality. The stronger theorem supplies the
existence of a matching feasible bound.
\begin{theorem}[Strong duality for linear programming]\label{thm:strong}
  If one member of a primal--dual LP pair has a finite optimal value and is
  feasible, both attain finite optima and those values are equal.
\end{theorem}
\begin{proof}[Proof through a certifying simplex basis]
  Convert the primal to full-row-rank standard form, removing redundant
  consistent equalities. Two-phase primal simplex with an anti-cycling rule
  terminates. Feasibility rules out failure in Phase I; a finite maximum rules
  out an improving unbounded ray. Its final basis has
  $\bar b\geq0$ and $q\geq0$. Define $y=A_B^{-T}c_B$.
  Then $A^Ty-c=q\geq0$, so $y$ is feasible for the equality-form dual, and
  \[
    b^Ty=c_B^TA_B^{-1}b=c_B^Tx_B=c^Tx.
  \]
  Weak duality certifies both optima. The equivalent transformations recover a
  dual certificate for the original form. Applying the same argument with the
  roles exchanged proves the assertion when the dual optimum is given first.
\end{proof}
An arbitrary basis representing a degenerate optimal point need not have
$q\geq0$. Use a \emph{certifying} optimal basis when extracting $y$ this way.

\begin{figure}[H]
  \centering
  \begin{tikzpicture}[every node/.style={font=\small}]
    \draw[model axis] (0,0)--(12,0) node[right] {objective value};
    \draw[modelblue,very thick] (1.3,-.12)--(1.3,.12);
    \draw[modelorange,very thick] (10.5,-.12)--(10.5,.12);
    \draw[black,thick] (6,-.16)--(6,.16);
    \node[below,align=center] at (1.3,-.15) {$c^Tx$\\feasible lower bound};
    \node[below,align=center] at (10.5,-.15) {$b^Ty$\\feasible upper bound};
    \node[above] at (6,.2) {$z^*=w^*$};
    \draw[model flow] (1.5,.65)--(5.7,.65);
    \draw[modelorange,thick,-{Stealth}] (10.3,.65)--(6.3,.65);
    \draw[<->] (1.3,1.25)--node[above] {duality gap $b^Ty-c^Tx\geq0$} (10.5,1.25);
  \end{tikzpicture}
  \caption{For a finite maximization LP, feasible primal values bound the optimum
    from below and feasible dual values bound it from above. Equality closes the gap.}
\end{figure}

\subsection{The outcome table}\label{sec:outcomes}
``Unbounded'' below means \emph{feasible with an unbounded objective}, not an
unbounded feasible set. The dual's unbounded direction is toward $-\infty$.
\begin{center}
  \begin{tabular}{lccc}\toprule
    Primal $\backslash$ dual & Finite optimum & Unbounded below & Infeasible\\\midrule
    Finite optimum & Possible; equal values & Impossible & Impossible\\
    Unbounded above & Impossible & Impossible & Possible\\
    Infeasible & Impossible & Possible & Possible\\\bottomrule
  \end{tabular}
\end{center}
An improving primal ray contradicts any finite feasible dual upper bound.
So primal unboundedness implies dual infeasibility. The converse does not
follow: both problems can be infeasible.
\begin{center}
  \begin{tabular}{>{\raggedright\arraybackslash}p{3.4cm}>{\raggedright\arraybackslash}p{5.1cm}>{\raggedright\arraybackslash}p{6cm}}\toprule
    Case & Primal & Dual\\\midrule
    Finite / finite & $\max x$, $0\leq x\leq1$ & $\min y$, $y\geq1$\\
    Unbounded / infeasible & $\max x$, $0x\leq1$, $x\geq0$ & $\min y$, $0y\geq1$, $y\geq0$\\
    Infeasible / unbounded & $\max0x$, $x\leq-1$, $x\geq0$ & $\min(-y)$, $y\geq0$\\
    Infeasible / infeasible & $\max x$, $0x\leq-1$, $x\geq0$ & $\min(-y)$, $0y\geq1$, $y\geq0$\\\bottomrule
  \end{tabular}
\end{center}
The zero rows make the last examples deliberately simple; no full-rank
standard-form assumption is needed to state these duality outcomes.

\section{Reading a dual certificate from an optimal tableau}\label{sec:production}
Return to the production model
\[
  \max(3x_1+2x_2),\qquad
  \begin{cases}
    2x_1+x_2\leq10,   \\
    3x_1+3x_2\leq24,  \\
    \tfrac32x_1\leq6, \\
    x_1,x_2\geq0.
  \end{cases}
\]
Introduce slacks $x_3,x_4,x_5$ in that row order. The final primal tableau is
\begin{center}
  \begin{tabular}{c|rrrrr|r}\toprule
    Basic & $x_1$ & $x_2$ & $x_3$ & $x_4$ & $x_5$ & RHS\\\midrule
    $x_2$ & 0 & 1 & $-1$ & $2/3$ & 0 & 6\\
    $x_5$ & 0 & 0 & $-3/2$ & $1/2$ & 1 & 3\\
    $x_1$ & 1 & 0 & 1 & $-1/3$ & 0 & 2\\\midrule
    $q^T$ & 0 & 0 & 1 & $1/3$ & 0 & 18\\\bottomrule
  \end{tabular}
\end{center}
Its basis is $B=(2,5,1)$; $c_B=(2,0,3)^T$. Solve $A_B^Ty=c_B$, rather than
forming an explicit inverse:
\[
  A_B=\begin{pmatrix}1&0&2\\3&0&3\\0&1&3/2\end{pmatrix},\qquad
  \begin{cases}
    y_1+3y_2=2, \\ y_3=0,\\2y_1+3y_2+\tfrac32y_3=3.
  \end{cases}
\]
Thus $y=(1,1/3,0)^T$. It is feasible for the dual
\[
  \min(10y_1+24y_2+6y_3),\quad
  2y_1+3y_2+\tfrac32y_3\geq3,\quad y_1+3y_2\geq2,\quad y\geq0,
\]
and $b^Ty=10+8=18=c^Tx^*$ for $x^*=(2,6)^T$.

More generally, for $[A\mid I]x=b$, $c_{\rm std}=(c,0)$,
\begin{equation}\label{eq:indicator-prices}
  q_{1:n}=A^Ty-c,\qquad q_{n+1:n+m}=y.
\end{equation}
The indicators under the primal slack columns are the dual variables. The
indicators under the original columns are the dual constraint surpluses.
This reading rule depends on the actual $+I$ slack convention; do not carry it
unchanged to surplus columns or rows multiplied by $-1$.

\begin{figure}[H]
  \centering
  \begin{tikzpicture}[every node/.style={font=\small}]
    \begin{scope}[x=.95cm,y=.43cm]
      \filldraw[model region] (0,0)--(4,0)--(4,2)--(2,6)--(0,8)--cycle;
      \draw[model axis] (0,0)--(5,0) node[right] {$x_1$};
      \draw[model axis] (0,0)--(0,9) node[above] {$x_2$};
      \draw[model objective] (0,9)--(5,1.5);
      \node[model point,label=right:{$(2,6)$}] at (2,6) {};
      \node at (2,-1.4) {Primal: $\max(3x_1+2x_2)$};
    \end{scope}
    \begin{scope}[xshift=7cm,x=1.9cm,y=1.9cm]
      \filldraw[model region] (0,1)--(1,1/3)--(2,0)--(2.3,0)--(2.3,1.8)--(0,1.8)--cycle;
      \draw[model axis] (0,0)--(2.6,0) node[right] {$y_1$};
      \draw[model axis] (0,0)--(0,2) node[above] {$y_2$};
      \draw[model objective] (0,.75)--(1.8,0);
      \node[model point,label=below left:{$(1,1/3)$}] at (1,1/3) {};
      \node at (1.1,-.32) {Dual slice: $y_3=0$};
      \node[align=center] at (1.1,1.35) {feasible region\\continues upward/right};
    \end{scope}
  \end{tikzpicture}
  \caption{The primal optimum and the dual certificate both have objective 18.
    The right drawing is the explicitly stated slice $y_3=0$, not a projection of
    the complete three-dimensional dual region. Dashed lines are objective level sets.}
\end{figure}

\subsection{Resource prices and their range of validity}\label{sec:prices}
The identity
\[
  3x_1+2x_2=(2x_1+x_2)+\tfrac13(3x_1+3x_2)
\]
prices the first two resources at $1$ and $1/3$ PLN per resource unit; the third
has price zero at this basis. For a change $b\mapsto b+\delta$, with $A,c$
fixed, the same basis remains optimal whenever $A_B^{-1}(b+\delta)\geq0$.
In that range,
\[
  z^*(b+\delta)=18+y^T\delta.
\]
This equality has a feasibility condition, not global validity for arbitrary
resource changes. For example, changing only the first RHS by $t$ gives
\[
  (x_2,x_5,x_1)=(6-t,\ 3-\tfrac32t,\ 2+t).
\]
The basis stays feasible exactly for $-2\leq t\leq2$, and revenue is $18+t$.
\begin{center}
  \begin{tabular}{crrrr}\toprule
    $t$ & $x_1$ & $x_2$ & Third slack & Revenue (PLN)\\\midrule
    $-2$ & 0 & 8 & 6 & 16\\
    0 & 2 & 6 & 3 & 18\\
    2 & 4 & 4 & 0 & 20\\\bottomrule
  \end{tabular}
\end{center}
At the endpoints the basis becomes degenerate; outside this interval, recompute
the optimum. A zero price and a tight constraint can coexist at degeneracy.

\section{Complementary slackness}\label{sec:complementarity}
For the max/min pair \eqref{eq:primal-dual}, take feasible $x,y$ and set
\[
  s=b-Ax\geq0,\qquad t=A^Ty-c\geq0.
\]
Their duality gap splits into nonnegative products:
\begin{equation}\label{eq:gap}
  b^Ty-c^Tx=y^T(b-Ax)+x^T(A^Ty-c)=y^Ts+x^Tt.
\end{equation}
\begin{theorem}[Complementary slackness]
  A primal-feasible $x$ and dual-feasible $y$ are both optimal if and only if
  \[
    y_is_i=0\quad(i=1,\ldots,m),\qquad x_jt_j=0\quad(j=1,\ldots,n).
  \]
\end{theorem}
\begin{proof}
  For optimal feasible points, strong duality gives zero gap. Since every summand
  in \eqref{eq:gap} is nonnegative, each product must vanish. Conversely, vanishing
  products give equal feasible objectives; weak duality certifies optimality.
\end{proof}
The points do \emph{not} need to be basic. Feasibility is essential: zero products
alone cannot certify an inadmissible point.
\begin{center}
  \begin{tabular}{ll}\toprule
    Positive quantity & Forced zero quantity\\\midrule
    Primal slack $s_i>0$ & Corresponding price $y_i=0$\\
    Price $y_i>0$ & Corresponding slack $s_i=0$\\
    Primal variable $x_j>0$ & Dual surplus $t_j=0$\\
    Dual surplus $t_j>0$ & Primal variable $x_j=0$\\\bottomrule
  \end{tabular}
\end{center}
These are implications, not equivalences. Both members of a complementary pair
can equal zero. In the production example, $s=(0,0,3)$, $y=(1,1/3,0)$,
$x=(2,6)$, and $t=(0,0)$, so every product vanishes.

\begin{figure}[H]
  \centering
  \begin{tikzpicture}[every node/.style={font=\small},
      box/.style={draw=modelblue,rounded corners,minimum width=2.4cm,minimum height=.8cm}]
    \node[box] (y) at (0,0) {price $y_i$};
    \node[box] (s) at (4.8,0) {resource slack $s_i$};
    \node[box] (x) at (0,-1.5) {quantity $x_j$};
    \node[box] (t) at (4.8,-1.5) {dual surplus $t_j$};
    \draw[modelblue,thick] (y)--node[above] {$y_is_i=0$} (s);
    \draw[modelblue,thick] (x)--node[above] {$x_jt_j=0$} (t);
    \node[align=left] at (9, -.75) {All four quantities $\geq0$\\Each pair: at least one zero\\Both zeros are allowed};
  \end{tikzpicture}
  \caption{Complementarity pairs each resource with its price, and each decision
    variable with the gap in its dual constraint.}
\end{figure}

\subsection{Certifying a supplied candidate without running simplex}\label{sec:candidate}
\paragraph{Task statement.}
\foreignlanguage{polish}{Sprawdzić, czy $x=[2,1,1,0,0]^T$ jest rozwiązaniem
optymalnym.}
Consider
\[
  \begin{aligned}
    \min\quad &x_1+4x_2+2x_3+2x_4+8x_5,\\
    &x_1+2x_2+x_4\geq4,\\
    &x_2+x_3-x_4\geq2,\\
    &x_2+x_5\geq1,\qquad x_1,\ldots,x_5\geq0.
  \end{aligned}
\]
For a min problem $Ax\geq b$, the dual is
\[
  \max b^Ty,\quad A^Ty\leq c,\quad y\geq0.
\]
Writing it explicitly gives
\[
  \begin{aligned}
    \max\quad &4y_1+2y_2+y_3,\\
    &y_1\leq1,\qquad 2y_1+y_2+y_3\leq4,\qquad y_2\leq2,\\
    &y_1-y_2\leq2,\qquad y_3\leq8,\qquad y_1,y_2,y_3\geq0.
  \end{aligned}
\]
The candidate is nonnegative and has $Ax=(4,2,1)^T=b$. Since its first three
coordinates are positive, complementary slackness forces their dual constraints
to be tight:
\[
  y_1=1,\quad 2y_1+y_2+y_3=4,\quad y_2=2.
\]
Hence $y=(1,2,0)^T$. The remaining inequalities give $-1\leq2$ and $0\leq8$,
so it is dual feasible. The min-form gap uses $s=Ax-b$ and $t=c-A^Ty$:
\[
  c^Tx-b^Ty=y^Ts+x^Tt.
\]
\begin{center}
  \begin{tabular}{crrr}\toprule
    $j$ & $x_j$ & $t_j=c_j-a_j^Ty$ & $x_jt_j$\\\midrule
    1 & 2 & 0 & 0\\
    2 & 1 & 0 & 0\\
    3 & 1 & 0 & 0\\
    4 & 0 & 3 & 0\\
    5 & 0 & 8 & 0\\\bottomrule
  \end{tabular}
\end{center}
All three primal surpluses are zero, so $y_is_i=0$ also holds. The objectives
match: $c^Tx=2+4+2=8=4+4=b^Ty$. Thus both candidates are optimal.
Moreover, positive dual surpluses force $x_4=x_5=0$, and $y_1,y_2>0$ force the
first two primal rows tight. Every primal optimum therefore has
\begin{equation}\label{eq:optimal-segment}
  x(t)=(4-2t,\ t,\ 2-t,\ 0,\ 0)^T,\qquad 1\leq t\leq2.
\end{equation}
The supplied candidate is the endpoint $t=1$, not the only optimizer.
The third row gives $t\geq1$; nonnegativity gives $t\leq2$.
The dual certificate is unique: any dual optimum must be complementary to the
supplied candidate, whose first three positive coordinates force $y=(1,2,0)$.

\begin{figure}[H]
  \centering
  \begin{tikzpicture}[every node/.style={font=\small}]
    \draw[model axis] (0,0)--(9,0) node[right] {$t=x_2$};
    \draw[modelblue,very thick] (2,0)--(7,0);
    \node[model point,label=below:{$t=1$: $(2,1,1,0,0)$}] at (2,0) {};
    \node[model point,label=below:{$t=2$: $(0,2,0,0,0)$}] at (7,0) {};
    \node[above,align=center] at (4.5,.5) {Entire segment optimal\\$c^Tx(t)=8$};
  \end{tikzpicture}
  \caption{The parameter line represents a one-dimensional optimal face in the
    five-dimensional decision space. Interior points of this segment are also optimal.}
\end{figure}

\section{The dual simplex method}\label{sec:dual-simplex}
For the standard max problem, a \emph{dual feasible basis} satisfies $q\geq0$;
its basic solution may have negative coordinates. The associated
$y=A_B^{-T}c_B$ satisfies $A^Ty\geq c$, and $b^Ty=z_0$, even if the primal
basic vector is infeasible. A negative $z_0$ or $\bar b$ is not by itself an error.
If also $\bar b\geq0$, primal and dual feasibility and equal values certify
optimality.
\begin{center}
  \begin{tabular}{lll}\toprule
    Feature & Primal simplex & Dual simplex\\\midrule
    Maintained condition & $\bar b\geq0$ & $q\geq0$\\
    Condition to repair & Some $q_k<0$ & Some $\bar b_r<0$\\
    First selection & Entering column & Leaving row\\
    Eligible pivot sign & $\bar a_{rk}>0$ & $\bar a_{rk}<0$\\
    Objective $z_0$ for max & Nondecreasing feasible value & Nonincreasing dual upper bound\\
    Failure certificate & Improving primal ray & Contradictory primal row\\\bottomrule
  \end{tabular}
\end{center}
The dual method is particularly convenient after changing $b$: indicators
depend on $A,c,B$, so they are unchanged, while $\bar b$ may lose nonnegativity.
It also naturally starts some min problems with nonnegative cost coefficients.

\subsection{The reversed ratio test}\label{sec:dual-ratio}
Choose a row $r$ with $\bar b_r<0$, for instance the most negative RHS.
Its equation is
\[
  x_{j_r}+\sum_{j\in N}\bar a_{rj}x_j=\bar b_r<0.
\]
If every nonbasic coefficient in that row is nonnegative, the left side cannot
be negative for any nonnegative vector. The primal is infeasible.
Otherwise a pivot can use a column with $\bar a_{rk}<0$.
Define $\alpha=q_k/(-\bar a_{rk})\geq0$. The new indicators are
\[
  q'_j=q_j+\alpha\bar a_{rj}.
\]
For $\bar a_{rj}\geq0$, nonnegativity is automatic. For negative coefficients,
we require $\alpha\leq q_j/(-\bar a_{rj})$. Therefore choose
\begin{equation}\label{eq:dual-ratio}
  k\in\operatorname*{arg\,min}_{j:\bar a_{rj}<0}
  \frac{q_j}{-\bar a_{rj}}
  \quad\Longleftrightarrow\quad
  k\in\operatorname*{arg\,max}_{j:\bar a_{rj}<0}
  \frac{q_j}{\bar a_{rj}}.
\end{equation}
This preserves dual feasibility. Since basic columns are unit columns, eligible
columns are automatically nonbasic. The new entering basic value
$\bar b_r/\bar a_{rk}$ is positive, but other basic values may still be negative.
Continue until all are nonnegative.

The objective update is
\[
  z'_0=z_0+\alpha\bar b_r\leq z_0.
\]
For a maximization problem this decreases a feasible \emph{dual upper bound},
not a feasible primal objective. If $q_k=0$, the dual step is degenerate and
this bound need not strictly decrease.

\begin{figure}[H]
  \centering
  \begin{tikzpicture}[x=1.1cm,y=.8cm,every node/.style={font=\small}]
    \draw[model axis] (0,0)--(4.8,0) node[right] {$\alpha$};
    \draw[model axis] (0,0)--(0,4.5) node[above] {indicator};
    \draw[modelblue,thick] (0,2)--(4,0);
    \node[modelblue,above] at (3.4,.55) {$q_1+\alpha v_1$};
    \draw[modelorange,thick] (0,3)--(1.5,0);
    \node[modelorange,above right] at (.35,2.3) {$q_2+\alpha v_2$};
    \draw[black,thick] (0,1)--(4,3);
    \node[right] at (4,3) {$q_3+\alpha v_3$};
    \draw[dashed] (1.5,0)--(1.5,3.5);
    \node[below,align=center] at (1.5,0) {$\alpha^*=3/2$\\column 2 enters};
  \end{tikzpicture}
  \caption{Added illustration: $q=(2,3,1)$ and pivot-row coefficients
    $v=(-1/2,-2,1/2)$. The first indicator to reach zero limits the dual step:
    $\min(2/(1/2),3/2)=3/2$. The positive-coefficient column imposes no upper limit.}
\end{figure}

\subsection{Algorithm and infeasibility certificate}\label{sec:dual-algorithm}
\begin{enumerate}[leftmargin=*,itemsep=2pt]
  \item Start with a nonsingular basis having $q\geq0$.
  \item If $\bar b\geq0$, return the optimal primal basic vector and its dual certificate.
  \item Choose a row $r$ with negative RHS.
  \item If that row has no negative nonbasic coefficient, certify infeasibility.
  \item Otherwise use \eqref{eq:dual-ratio} and pivot on its negative element.
  \item Recompute or update the tableau and return to step 2.
\end{enumerate}
A smallest-column tie break makes a worked path reproducible, but the simple
``most negative RHS'' rule alone is not an anti-cycling theorem. General dual
simplex needs a proven anti-cycling variant when degeneracy can occur.

To express failure without a tableau, put $h=A_B^{-T}e_r$. Then the
contradictory row satisfies
\[
  h^TA=\bar A_{r:}\geq0,\qquad h^Tb=\bar b_r<0.
\]
Any $x\geq0$ with $Ax=b$ would give
$h^Tb=h^TAx\geq0$, a contradiction. This is an explicit infeasibility
certificate. Starting at any dual-feasible $y$, the ray $y+\lambda h$ remains
dual feasible and its objective tends to $-\infty$.

For example, $x_1+x_2=-1$, $x\geq0$, $\max0$ has the dual-feasible basis
$B=(1)$ with $q=(0,0)$ and RHS $-1$. Its row has no negative coefficient.
The certificate is $h=1$, and the dual $\min(-y)$, $y\geq0$ is unbounded below.
This single row simultaneously explains the two outcomes.

\section{A complete dual-simplex solve}\label{sec:dual-example}
Consider
\[
  \begin{aligned}
    \min\quad &x_1+x_2+x_3,\\
    &x_1+2x_2+x_3\geq2,\\
    &2x_1+x_2-x_3\geq10,\qquad x_1,x_2,x_3\geq0.
  \end{aligned}
\]
For a general min problem $\min c^Tx$, $Ax\geq b$, $x\geq0$ with $c\geq0$,
negate the objective and rows and add surplus variables:
\[
  \max(-c^Tx),\qquad -Ax+s=-b,\qquad x,s\geq0.
\]
The surplus basis is $I$, with $x=0$, $s=-b$ and indicators $(c,0)\geq0$.
It is dual feasible regardless of the signs of $b$, and primal feasible only
when $-b\geq0$. Thus no artificial variables are needed for this start.

Here $x_4,x_5$ denote the nonnegative surpluses of the original inequalities:
\[
  -x_1-2x_2-x_3+x_4=-2,\qquad
  -2x_1-x_2+x_3+x_5=-10.
\]
The start $B=(4,5)$ is primal infeasible but dual feasible:
\begin{center}
  \begin{tabular}{c|rrrrr|r}\toprule
    Basic & $x_1$ & $x_2$ & $x_3$ & $x_4$ & $x_5$ & RHS\\\midrule
    $x_4$ & $-1$ & $-2$ & $-1$ & 1 & 0 & $-2$\\
    $x_5$ & $-2$ & $-1$ & 1 & 0 & 1 & $-10$\\\midrule
    $q^T$ & 1 & 1 & 1 & 0 & 0 & 0\\\bottomrule
  \end{tabular}
\end{center}
The most negative RHS is in row 2, so $x_5$ leaves. Only columns 1 and 2
have negative entries there:
\begin{center}
  \begin{tabular}{crrr}\toprule
    Column & $q_j$ & $\bar a_{2j}$ & $q_j/(-\bar a_{2j})$\\\midrule
    1 & 1 & $-2$ & $1/2$\\
    2 & 1 & $-1$ & 1\\\bottomrule
  \end{tabular}
\end{center}
Choose $x_1$ to enter, and pivot on $-2$. Divide row 2 by $-2$, add that
normalized row to row 1, and subtract it from the indicator row:
\begin{center}
  \begin{tabular}{c|rrrrr|r}\toprule
    Basic & $x_1$ & $x_2$ & $x_3$ & $x_4$ & $x_5$ & RHS\\\midrule
    $x_4$ & 0 & $-3/2$ & $-3/2$ & 1 & $-1/2$ & 3\\
    $x_1$ & 1 & $1/2$ & $-1/2$ & 0 & $-1/2$ & 5\\\midrule
    $q^T$ & 0 & $1/2$ & $3/2$ & 0 & $1/2$ & $-5$\\\bottomrule
  \end{tabular}
\end{center}
Both RHS entries are now nonnegative and all indicators remain nonnegative.
Thus $(x_1,x_2,x_3)=(5,0,0)$, surpluses $(3,0)$, and the \emph{original
  minimum} is 5. The transformed maximum is $-5$.

The original min-form dual is
\[
  \max(2y_1+10y_2),\quad
  y_1+2y_2\leq1,\quad2y_1+y_2\leq1,\quad y_1-y_2\leq1,\quad y\geq0.
\]
It has certificate $y=(0,1/2)$, objective 5, and column slacks $(0,1/2,3/2)$.
Equivalently, solve $A_B^Ty=c_B$ for the transformed max to get the same
nonnegative $y$: its surplus columns form $+I$ after row negation.
The original primal inequalities directly imply the bound
\begin{equation}\label{eq:dual-example-certificate}
  x_1+x_2+x_3
  =5+\tfrac12x_2+\tfrac32x_3+\tfrac12x_5\geq5.
\end{equation}
Equality forces $x_2=x_3=x_5=0$, then $x_1=5$; the optimum is unique.

\begin{figure}[H]
  \centering
  \begin{tikzpicture}[every node/.style={font=\small},
      box/.style={draw=modelblue,rounded corners,align=center,minimum width=4.8cm,minimum height=1.2cm}]
    \node[box] (before) at (0,0) {$B=(4,5)$\\$\bar b=(-2,-10)$, $q\geq0$\\dual bound $0$ for the max};
    \node[box] (after) at (7,0) {$B=(4,1)$\\$\bar b=(3,5)$, $q\geq0$\\dual bound $-5$: attained};
    \draw[model flow] (before)--node[above] {pivot $-2$} (after);
    \node[align=center] at (3.5,-1.5) {Original min: lower bound rises from $0$ to $5$\\Primal feasibility is repaired in one pivot};
  \end{tikzpicture}
  \caption{Dual simplex keeps a valid dual certificate throughout. The starting
    basic vector is not an admissible solution of the original minimization problem.}
\end{figure}

\section{MATLAB checks and study questions}\label{sec:matlab}
The companion uses base MATLAB and the same matrices as the hand examples:
\begin{lstlisting}
run('matlab/exercises/lecture_04_checks.m')
\end{lstlisting}
Basis computations use linear solves, not an explicit inverse:
\begin{lstlisting}
AB = A(:, B);
Abar = AB \ A;
bbar = AB \ b;
y = AB' \ c(B);
q = A' * y - c;
x = zeros(size(A, 2), 1);
x(B) = bbar;
\end{lstlisting}
The assertions check specified small examples, not a general-purpose LP solver.
They compare known paths, feasible points, certificates, and sign conditions
with absolute tolerance $10^{-10}$. Lexicographic comparisons use the same
tolerance for these moderate-size examples; this is not exact symbolic arithmetic
or a universal anti-cycling guarantee for ill-conditioned data.

\begin{center}
  \begin{tabular}{>{\raggedright\arraybackslash}p{5.6cm}>{\raggedright\arraybackslash}p{9cm}}\toprule
    Check & Expected hand result\\\midrule
    Cycling example & Six zero steps; starting basis repeats\\
    Bland / lexicographic paths & Six / two pivots; both attain $5/4$\\
    Three-dimensional distorted cube & Eight vertices, seven pivots, maximum 125\\
    Production dual and RHS changes & $y=(1,1/3,0)$; $18+t$ for $-2\leq t\leq2$\\
    Complementary-slackness candidate & Value 8; entire optimal segment $1\leq t\leq2$\\
    Dual-simplex solve & One pivot; minimum 5; $y=(0,1/2)$\\
    Contradictory tableau row & $h^TA\geq0$, $h^Tb<0$: infeasible\\\bottomrule
  \end{tabular}
\end{center}
\subsection{Common mistakes}
\begin{center}
  \begin{tabular}{>{\raggedright\arraybackslash}p{6cm}>{\raggedright\arraybackslash}p{8.6cm}}\toprule
    Mistake & Repair\\\midrule
    Calling every zero coordinate degenerate & Check the basic coordinates\\
    Equating degeneracy with cycling & Distinguish a zero step from a repeated basis\\
    Using the max sign table for a min primal & Reverse the appropriate signs, or negate the objective first\\
    Treating a tight row as a positive-price guarantee & Complementarity gives implications; both quantities can be zero\\
    Checking only equality of objective values & Verify feasibility on both sides\\
    Reading all indicator rows as resource prices & Identify the actual original slack columns and row signs\\
    Using the primal ratio test in dual simplex & Choose a negative RHS row first, then minimize $q_j/(-\bar a_{rj})$\\
    Reporting transformed value $-5$ as the minimum & Undo objective negation: the minimum is 5\\\bottomrule
  \end{tabular}
\end{center}
\Needspace{20\baselineskip}
\subsection{Questions with answers for checking}
\begin{enumerate}[leftmargin=*,itemsep=2pt]
  \item In the cycling example's fourth tableau, why do Dantzig and Bland choose
        different entering variables?
  \item Why does the perturbation rule select row 2 in the first ratio tie?
  \item For $\max2x$, $x=3$, $x\geq0$, construct the dual and its optimum.
  \item Does a zero resource price prove a resource is unused?
  \item Derive every optimal primal point for the complementary-slackness task.
  \item In the dual-simplex example, what happens if column 2 is chosen at the
        first pivot instead of column 1?
\end{enumerate}
\paragraph{Answers.}
(1) Dantzig chooses $q_7=-3$, while Bland chooses the smaller eligible index 1
with $q_1=-2$.
(2) The normalized tuples start $(0,4,0,\ldots)$ and $(0,0,2,\ldots)$.
(3) $\min3y$, $y\geq2$, $y$ free; optimum $y=2$, value 6.
(4) No: $y_i=0$ permits either zero or positive slack.
(5) Equation~\eqref{eq:optimal-segment} follows from $x_4=x_5=0$ and tight
first two rows; all its points have value 8.
(6) Its ratio 1 exceeds the allowed $1/2$; the updated $q_1=1-2=-1$ loses
dual feasibility. An arbitrary negative pivot is insufficient.

\clearpage
\section{Sources and editorial remarks}\label{sec:sources}
\begingroup\small
\noindent\textbf{Primary source:} E. Pawelec, \texttt{s4\_w4.pdf},
PDF pages 1--12, printed slides 1--47 (slide 47 is the closing slide).\\
\textbf{Source date:} The folder is labelled 2020; PDF creation metadata is
2019-10-27, not an established delivery date.\\
\textbf{Current course alignment:} Supplied sheet~1, groups~2--5, requires
dual construction, complementarity and dual simplex. These topics are not
developed in the supplied \texttt{ProgMat\_w1.pdf}; its numbering does not
establish when they were taught. Cycling and complexity remain historical
enrichment.\\
\textbf{Editorial scope:} Complete mathematical lecture content, reorganized
in English, with hand proofs, worked paths, certificates, study questions, and
editable TikZ illustrations. All source pages were reviewed; dense cycling,
perturbation, and final pivot tableaus were also inspected at enlarged scale.

\begin{center}
  \begin{tabular}{lll}\toprule
    PDF pages & Printed slides & Developed material\\\midrule
    1 & 1--4 & Section~\ref{sec:notation}: standard form and pivot equations\\
    2 & 5--8 & Section~\ref{sec:degeneracy}: degeneracy and full cycling tableaus\\
    3--4 & 9--14 & Sections~\ref{sec:perturbation}, \ref{sec:bland}: perturbation and Bland\\
    4--5 & 15--18 & Section~\ref{sec:complexity}: complexity and geometric routes\\
    5--6 & 19--23 & Section~\ref{sec:dual-rules}: dual construction and mixed signs\\
    6--7 & 24--28 & Section~\ref{sec:duality}: weak/strong duality and outcomes\\
    7--8 & 29--31 & Section~\ref{sec:production}: extracting the production dual\\
    8--9 & 32--36 & Section~\ref{sec:complementarity}: optimality via complementarity\\
    10--11 & 37--44 & Section~\ref{sec:dual-simplex}: dual simplex and its ratio rule\\
    12 & 45--46 & Section~\ref{sec:dual-example}: initialization and complete solve\\\bottomrule
  \end{tabular}
\end{center}
The Polish practice prompt in Section~\ref{sec:bland} is from slide 14.
The candidate task in Section~\ref{sec:candidate} preserves the mathematical
question on slide 34 with spelling normalized; its data are from slides 34--36.
The full Bland path, the $p=5$ three-dimensional vertex list, mixed-example
feasible pair, resource-price interval, optimal segment, and explicit row
infeasibility certificate are added derivations. Drawings explain the source
ideas in the current style; the ratio-indicator graph uses added data.

\paragraph{Qualifications.}
Section~\ref{sec:degeneracy} distinguishes zero steps from positive tied steps,
and allows a negative-indicator basis at a degenerate optimal point.
Section~\ref{sec:perturbation} makes the initial feasibility requirement explicit.
Section~\ref{sec:complexity} attaches the distorted-cube route to its pivot rule
and does not adopt slide 16's unspecified average-complexity expression.
Section~\ref{sec:duality} distinguishes the equal-value certificate on slide 25
from strong duality's existence statement. A certifying basis is needed to
extract dual multipliers from an optimal point. Complementarity applies to
feasible points, including nonbasic ones. Section~\ref{sec:dual-algorithm}
qualifies termination under dual degeneracy.

\paragraph{Supplementary mathematical checks.}
MIT OCW 15.081J/6.251J, Fall 2009:
\href{https://ocw.mit.edu/courses/6-251j-introduction-to-mathematical-programming-fall-2009/2ef2f1dd7045b5f29e5faea299fb0798_MIT6_251JF09_lec06.pdf}{Lecture 6, The Simplex Method II},
PDF pp.~5--8, for symbolic perturbation and anti-cycling;
\href{https://ocw.mit.edu/courses/6-251j-introduction-to-mathematical-programming-fall-2009/496d9a385fd1d7e1aa82837d91adb23e_MIT6_251JF09_lec08.pdf}{Lecture 8, Duality Theory I},
PDF pp.~2--5, for dual signs, the strong-duality basis proof, and complementarity;
\href{https://ocw.mit.edu/courses/6-251j-introduction-to-mathematical-programming-fall-2009/b93d1665d7a8f05ea626b8f04c7ad66c_MIT6_251JF09_lec09.pdf}{Lecture 9, Duality Theory II},
PDF pp.~4--5, for dual-simplex invariants and the ratio test.
Titles here follow the PDF interiors, which differ from some catalogue labels.
These references were consulted through the web tool; shell downloads failed
at DNS resolution. Intended storage is outside this repository, under
\texttt{../sources/mit/}; the paired source inventory records URLs and status.
They do not establish current course coverage.

\paragraph{Verification.}
Exact rational calculations independently recomputed the cycling, Bland,
lexicographic, distorted-cube, and dual-simplex paths and their certificates.
The installed R2026b static analyzer reported no diagnostics. On 2026-10-04,
the user ran the MATLAB companion in a normal terminal; all numerical assertions
passed and MATLAB reported version 26.2.0.3386108 (R2026b). The supplied output
agrees with Section~\ref{sec:matlab}. Execution evidence is recorded in
\path{docs/matlab-verification.md}; Polish terms are in \path{docs/glossary.md}.
\endgroup
\end{document}
